\chapter{Background} \label{chap:background}
\section{Network sockets}
Socket programming today exists in several forms and can
in several forms, and can be used for both inter-process communication
and network communication. When it comes to network programming,
the network socket in its bost basic definition is an abstraction
for a network connection, through which an application can send
and receive data \cite[p. 8]{donahooTCPIPSockets2009}. 
First mentioned in 1970 in relation to ARPANET, the
simplicity of sockets and their ubiquity in all UNIX distributions
since version 4.2 of the Berkeley Software Distribution have
led to them becoming the de facto standard for low-level network
programming \cite{crockerNewHostHostProtocol1970}\cite[p. 502]{tanenbaumComputerNetworks2011}.

Modern sockets are standardized across UNIX-like systems through
the POSIX standard. Implementations of the C POSIX library define
the socket API primarily through the \texttt{sys/socket.h} header file
\cite{posix:2024:socket}.

When creating a socket on an UNIX-like system, the application
invokes the \texttt{socket} system call. This system call
\say{creates an endpoint for communication and returns a file 
descriptor that refers to that endpoint} \cite{Socket2LinuxManual}.
This system call takes three parameters, the domain, the type
and the protocol which will be used for the socket.
The domain specifies the type of
addressing which will be used, for example \texttt{AF\_UNIX}
for local sockets, \texttt{AF\_INET} for IPv4 and \texttt{AF\_INET6}
for IPv6. The type specifies the protocol which will be used over the socket, for
example \texttt{SOCK\_DGRAM} for datagrams and \texttt{SOCK\_STREAM}
for a reliable stream of bytes. The type usually corresponds to a
single transport layer protocol, usually UDP and TCP respectively.
The protocol parameter is used if there are multiple protocols
which support the requirements given by the domain and type.

The fact that the application has to specify both the domain
and the type of the socket means that both the
underlying addressing and the protocol are bound at compile time.
The underlying transport protocol is therefore tightly
coupled to the design of the application, and barring fairly
advanced logic when setting up
the necessary sockets, the application will be bound to a single 
protocol and addressing scheme. Updating
the underlying protocol therefore requires at least recompilation and 
probably a fairly large refactoring of the application. 
This naturally makes adopting new protocols difficult, even
if the new protocol is better suited to the application's requirements.

Another challenge with new transport-layer protocols apart
from TCP and UDP is that developing them has proven to be challenging, 
in part because of TCP and UDP's ubiquity. 

Network middleboxes have have been
optimized to support TCP and UDP, and may sometimes
even drop packets if they contain another transport layer protocol.
The resulting difficulty in innovating on the transport-layer
is known as network ossification. Several new protocols such
as UDP-lite and QUIC are built on top of UDP in order to combat
this ossification.

While using UDP as a substrate does present some overhead,
it is fairly minimal with only an 8 bytes.

Applications can themselves implement needed features
such as reliability and congestion control.
This way of designing new protocols has become so common that
it has prompted an official RFC from the IETF, describing
the best practices for implementing congestion control on top of UDP 
\cite{eggertUnicastUDPUsage2008}.

\section{Transport Services}
At its core, a protocol is an abstraction which provides a set of
services to an upper-layer protocol or application. In the context of
networking, these services encompass transporting data from one endpoint
to another, with the individual features of a transport service deciding the behavior of
this service, for example when experiencing packet loss, packet reordering
or other issues. For example, TCP provides a guarantee of reliable
delivery, while UDP provides no such guarantee.

At their core, transport layer protocols are chosen based
on the services provided to the application, not because of the
protocols themselves. If an application requires that the
transferred data is received as an in-order stream, then
any protocol offering that as a service to the application layer
could function as a transport layer protocol, and swapping the
underlying transport layer protocol should in theory be simple.

This concept of exactly which services each transport layer protocol
provides to the application layer is known as Transport Services.
Exactly which Transport Services are provided by which transport
protocols has been mapped by the IETF in RFC8095, as part of the
TAPS working group \cite{fairhurstServicesProvidedIETF2017}.

\section{What is TAPS}
The TAPS WG is an IETF working group which aims to solve the strong
coupling between application
level code and the underlying transport protocol through
focusing on Transport Services. This is intended to allow faster adaptation
of new protocols and abstract away network-layer details from the
application layer. From the application's point of view, this
means that it will be able to gain access to new protocols, such
as QUIC, without having to implement its own fallback logic
in cases where the new protocl is not supported by the remote.
It will also simplify the required networking logic by
tying it more closely to how modern applications actually interact the network.

Despite being implemented in user-space, TAPS conceptually resides
in the transport layer. TAPS provides convenience to the application
layer by expanding the responsiblity of the transport layer. Instead
of the application layer being responsible for resolving addresses and
for chosing the specific transport layer API to use,
TAPS provides a unified interface which
both handles simple interactions with the network for resolving hostnames
in addition to unifying the transport protocol APIs. This minimizes
how much the application layer has to adapt to the underlying transport
layer, and minimizes the coupling between the two layers.

The difference in API models for socket and for TAPS can be
seen in \cref{fig:socket-programming-api-model} and
\cref{fig:taps-api-model}.

\begin{figure}[tbp]
    \centering
    \begin{lstlisting}
        +-----------------------------------------------------+
        |                    Application                      |
        +-----------------------------------------------------+
                |                 |                  |
          +------------+     +------------+    +--------------+
          |  DNS Stub  |     | Stream API |    | Datagram API |
          |  Resolver  |     +------------+    +--------------+
          +------------+          |                  |
                            +---------------------------------+
                            |    TCP                UDP       |
                            |    Kernel Networking Stack      |
                            +---------------------------------+
                                            |
        +-----------------------------------------------------+
        |               Network-Layer Interface               |
        +-----------------------------------------------------+
    \end{lstlisting}
    \caption{API model for socket programming \cite{trammellAbstractApplicationProgramming2025}}
    \label{fig:socket-programming-api-model}
\end{figure}

\begin{figure}[tbp]
    \centering
    \begin{lstlisting}
        +-----------------------------------------------------+
        |                    Application                      |
        +-----------------------------------------------------+
                                  |
        +-----------------------------------------------------+
        |              Transport Services API                 |
        +-----------------------------------------------------+
                                  |
        +-----------------------------------------------------+
        |          Transport Services Implementation          |
        |  (Using DNS, UDP, TCP, SCTP, DCCP, TLS, QUIC, etc.) |
        +-----------------------------------------------------+
                                  |
        +-----------------------------------------------------+
        |               Network-Layer Interface               |
        +-----------------------------------------------------+
    \end{lstlisting}
    \caption{API model for TAPS \cite{trammellAbstractApplicationProgramming2025}}
    \label{fig:taps-api-model}
\end{figure}

\subsection{TAPS Architecture} \label{sec:tapsArch}
RFC 9621 describes the underlying architecture of and a set of requirements
for TAPS, a diagram of this architecture can be seen in \cref{fig:tapsArch} \cite{paulyArchitectureRequirementsTransport2025}.
A TAPS system consists of two primary parts, the Transport Services
API and the Transport System Implementation. The Transport Services
API is the public facing API used by the Application. This contains
concepts defined in RFC 9622 such as a Preconnection, a Connection
Messages or Events being transferred between TAPS and the Application.

The Transport System Implementation is the internal workings
of a TAPS system. This consists of services which are provided to the
public API, such as data transfer or Candidate Racing. The TAPS RFCs
do not prescribe an actual detailed implementation of the Transport
System Implementation, leaving the details up to each TAPS implementation
itself. RFC 9623 does however offer a set of suggestions to developers
of a TAPS API in order to be in line with the TAPS specification,
however these are not mandatory, and any TAPS system is therefore
free to implement its internal workings as it sees fit.

\begin{figure}[tbp]
    \centering
    \begin{lstlisting}
     +-----------------------------------------------------+
     |                    Application                      |
     +-+----------------+------^-------+--------^----------+
       |                |      |       |        |
     pre-               |     data     |      events
     establishment      |   transfer   |        |
       |        establishment  |   termination  |
       |                |      |       |        |
       |             +--v------v-------v+       |
     +-v-------------+   Connection(s)  +-------+----------+
     |  Transport    +--------+---------+                  |
     |  Services              |                            |
     |  API                   |  +-------------+           |
     +------------------------+--+  Framer(s)  |-----------+
                              |  +-------------+
     +------------------------|----------------------------+
     |  Transport             |                            |
     |  System                |        +-----------------+ |
     |  Implementation        |        |     Cached      | |
     |                        |        |      State      | |
     |  (Candidate Gathering) |        +-----------------+ |
     |                        |                            |
     |  (Candidate Racing)    |        +-----------------+ |
     |                        |        |     System      | |
     |                        |        |     Policy      | |
     |             +----------v-----+  +-----------------+ |
     |             |    Protocol    |                      |
     +-------------+    Stack(s)    +----------------------+
                   +-------+--------+
                           V
   +-----------------------------------------------------+
   |               Network-Layer Interface               |
   +-----------------------------------------------------+
    \end{lstlisting}
    \caption{TAPS architecture as described by RFC 9621 \cite{paulyArchitectureRequirementsTransport2025}.}
    \label{fig:tapsArch}
\end{figure}

\subsection{TAPS abstractions}
Abstracting away the underlying transport protocols is done by
presenting a set of abstractions to the application. The most
central of these abstractions is the Connection. A Connection
represents a \say{shared state of two or more endpoints that persists
across messages that are transmitted between these endpoints}
\cite{welzlUsageTransportFeatures2018}.
When working with Connection objects, TAPS provides five high
level \textit{groups of actions} that an application can perform:
\begin{itemize}
    \item \texttt{Preestablishment}: Pass properties to specify
        requirements and preferences for establishing a Connection.
    \item \texttt{Establishment}: Prepare a Connection object for data
        transfer.
    \item \texttt{Data transfer}: Sending and receiving messages
        over the Connection.
    \item \texttt{Event handling}: Notifications that an application
        can receive from a Connection, and the handling of them.
    \item \texttt{Termination}: How transmission is stopped and
        a Connection is torn down.
\end{itemize}

In addition, the lifeline of a Connection is divided into three distinct
phases, \textit{preestablishment}, the \textit{Established state}, and 
\textit{termination of a Connection} \cite{paulyArchitectureRequirementsTransport2025}.
Each of these phases are associated
with one or more of the \textit{groups of actions} that an application can perform,
as described above. 

After the desired Transport Properties and cryptographic parameters have
been set during the preestablishment phase, a connection can be established. 
When a Connection is in the established phase the Connection can be used to
transfer data. 

There are three ways of
establishing a connection: \textit{initiate}, \textit{listen} and
\textit{rendezvous} \cite{paulyArchitectureRequirementsTransport2025}.

\textit{Initiate} is used to establish a connection to a remote endpoint which
is assumed to be listening for incoming connections. How this connection
is established internally is dependent on the underlying transport
protocol, which has been selected based on the specified Transport Properties.
The Created Connection object is passed to the application through a callback.

\textit{Listen} is used to passively open a connection. This is done in order to
wait for connections from remote endpoints. Listen creates a Listener
object. When a remote endpoint connects to this listener, a \texttt{ConnectionReceived}
event occurs. This event contains a Connection object for the 
established connection.

\textit{Rendezvous} is used to establish a peer-to-peer connection between two endpoints
simultaneously.

After a Connection has been established, it can either
be closed or aborted, with closing being associated with normal termination,
and aborting being abnormal termination associated with errors.

One or more Connection can be organized into a Connection Group. THese
are used to group related connections together, which have been created
by cloning. Connections in the same connection group can be
multiplexed together if this is supported by the underlying protocol.

In addition to the Connection, TAPS also provides several other objects
abstracting key elements. For example Endpoint objects and Message objects.

\subsection{Framer}
A Framer is another key concept within TAPS. A Framer is responsible for briding the TAPS concept
of a Message, with the underlying protocols concept of either a data stream
or a datagram. In their simplest form, a Framer is responsible for encoding
or decoding outgoing and incoming data, as well as sending any data necessary
on initiation or closing of a connection. Because Framers may need to send
data upon the creation of a Connection, they must be added to a Preconnection,
before any action which may result in the creation of a Connection is taken.
The application can define custom framers and pass them to a preconnection
by implementing an interface provided by the TAPS system.

Because of this potential need to interact with the network, 
framers are also asynchronous. They are therefore callback-based and need
to notify their associated connection when the processing of a given message
has completed.

The concept of a framing protocol exists outside of TAPS. Because most
application-level protocols are both message-based and run over TCP, several
ways of converting TCP's byte-stream into usable messages has been implemented.
These largely fall into two categories, type-length-value (TLV) structures where
the length of the following data is described in a header-structure, and
delimeter-separated formats, where the beginning and end of a piece of
data is separated by delimeters. \cite{brunstromImplementingInterfacesTransport2025}
Examples of the two categories are HTTP/1.1 for a delimeter-separated format
and DHCP for TLV, specifying the length of the given hardware address, allowing
parsers which are unfamiliar with the type to skip it, since the length is known
\cite{dromsDynamicHostConfiguration1997}.

For a TAPS system, where the appliaction layer should not need to take into
account which application layer protocol it is running on top of, Framers become
useful for datagram or message-based protocols as well. By using Framers, the
application layer protocol can ensure a consistent reception of messages in
the desired format. Framers can also pack multiple messages into a single transport
datagram by using the same formatting as it would for a stream-based protocol.

Conceptually, a framer resides on a separate layer between a Connection and the
underlying Transport Protocol Stack as seen in \cref{fig:tpsStack}.

\begin{figure}[tbp]
  \begin{lstlisting}
     Initiate()   Send()   Receive()   Close()
       |          |         ^          |
       |          |         |          |
  +----v----------v---------+----------v-----+
  |                Connection                |
  +----+----------+---------^----------+-----+
       |          |         |          |
       |      +-----------------+      |
       |      |    Messages     |      |
       |      +-----------------+      |
       |          |         |          |
  +----v----------v---------+----------v-----+
  |                Framer(s)                 |
  +----+----------+---------^----------+-----+
       |          |         |          |
       |      +-----------------+      |
       |      |   Byte-stream   |      |
       |      +-----------------+      |
       |          |         |          |
  +----v----------v---------+----------v-----+
  |         Transport Protocol Stack         |
  +------------------------------------------+
  \end{lstlisting}
\caption{Protocol Stack Showing a Message Framer from RFC 9622 \cite{trammellAbstractApplicationProgramming2025}}
\label{fig:tpsStack}
\end{figure}

Framers can also be configured on a per-message basis using a MessageContext.
This MessageContext can for example be used to pass the type of a message
to a Framer if the framer needs it to follow a TLV format.

\subsection{Transport Properties}
Transport Properties are properties which are set in TAPS in order to either
influence which transport protocol should be used, or how the select transport
protocol should behave during the lifetime of the Connection.

\subsubsection{Selection Properties}
Selection properties are the subset of Transport Properties which
used to select the protocol stack
to use for a given connection when performing candidate gathering as described
in \cref{sec:candidate-gathering}.

Typical Selection Properties are specified as
a \textit{preference}, with the possible values being \textit{require},
\textit{prefer}, \textit{no preference}, \textit{avoid} and \textit{prohibit}.
\textit{Require} means that the Connection establishment will fail if the property cannot
be satisfied, while \textit{prefer} will allow the application to continue.
This distinction is mirrored in the difference between \textit{prohibit} and \textit{avoid}.
Not all Selection Properties are set as a preference,
instead specified for example as a boolean, an integer or an enum value.

An example of a selection property is Reliable Data Transfer.
This property is used to specify whether or not the application
should use a protocol which ensures that all data is received at 
a remote endpoint, in order, without loss or duplication. Setting
this to \textit{require} will force TAPS to use an underlying protocol
which guarantees this, such as TCP or QUIC. Setting this as
\textit{prohibit} will force TAPS to use a protocol which does not
guarantee this, such as UDP.

\subsubsection{Connection Properties}
Connection Properties are the subset of Transport Properties which
are used to control the behavior of the underlying protocol after
a Connection has been established.

An example of a connection property is Timeout for Aborting Connection.
This property is used to specify the how long to wait before deciding that an
active connection has failed when trying to reliably deliver data to the Remote Endpoint.
This naturally can only take effect if the underlying protocol supports reliability.

\subsection{Message Properties}
Message properties are settings set on an
individual message. These are used to set the behavior
of a single message when transferring data, but does
not affect any other messages transferred by the Connection.

An example of a Message Property is \mono{Safely Replayable}.
This property is used to let TAPS know that a message is idempotent
and does not pose a risk if it is replayed to the endpoint. This
is for example a requirement for QUIC 0-RTT, in order to mitigate
the risk of replay attacks.

\subsection{TAPS configurations explained}
As mentioned above, TAPS allows the application to control
both the behavior and selection of the underlying transport
protocol through the use of \textit{Transport Properties}.

In addition to the basic Transport Properties such as reliability
or timeout, TAPS also exposes several more advanced
Transport Properties and these are explained below.

The \texttt{Multistream Connections in Group} Selection Property
allows the application to select whether the
application would prefer multiple Connections within a
Connection Group to be provided by streams of a single
underlying transport connection where possible
\cite{paulyArchitectureRequirementsTransport2025}.
Multistreaming in this context is the act of providing multiple
streams of data over a single underlying connection. This is
advantageous as it allows for more efficient use of the network,
since the connection will better saturate the available
bandwidth. This avoids long convergence times for the connection
and allows fresh connection to utilize the large congestion
window already achieved, without having to wait for slow start
to ramp up. This is especially useful for applications with many
small streams of data, such as web browsers.

The \mono{Full Checksum Coverage on Sending} and equivalent
Selection Property for receiving,
\mono{Full Checksum Coverage on Receiving}, specify
the application's need for protection against corruption for all
data transmitted or received on a given connection. This is used
to specify how many bytes of the message which needs to be
protected by a checksum. This checksum is used to protect the
data from corruption caused by lower-layer protocols. The
setting \mono{Full coverage} means that the entire message is
protected, while \mono{0} specifies no checksum is required.
If any other value than \mono{Full coverage} is specified, the
connection may still fall back to \mono{Full coverage}.

The \mono{Provisioning Domain Instance or Type} selection
property allows the application
to control path selection by selecting for specific Provisioning
Domain or categories of domains
\cite{paulyArchitectureRequirementsTransport2025}. Provisioning
Domains define consistent sets of network properties that might
be more specific than network interfaces
\cite{paulyArchitectureRequirementsTransport2025}.

The \texttt{Use Temporary Local Address} selection property
allows the application to specify
whether or not to use temporary addresses which are a feature in
IPv6, used for greater privacy
\cite{gontTemporaryAddressExtensions2021},\cite{paulyArchitectureRequirementsTransport2025}.

The \texttt{Notification of ICMP soft error message arrival}
selection property lets an application specify whether or 
not it should be notified when an ICMP
error, which does not force termination of the connection,
arrives \cite{paulyArchitectureRequirementsTransport2025}.

In addition to specifying which underlying protocol should be
used through selection properties, the application using TAPS
can manage the connections
themselves through connection properties. These can be specific
for a given transport protocol
or generic and be available for all transport protocols.

The \texttt{Connection Group Transmission Scheduler} connection 
property decides which scheduler is used within a given
connection group, to divide the available capacity between the
connections. \cite{paulyArchitectureRequirementsTransport2025}.

The \texttt{Capacity Profile} connection property is used to
select a capacity profile, which will be a tradeoff between
delay, delay variation, and efficient use of the available
capacity based on the capacity profile specified
\cite{paulyArchitectureRequirementsTransport2025}.

The \texttt{Isolate Session} connection property set in order to
have new connections use as little cacheed information as
possible from previous connections not in the same connection
group \cite{paulyArchitectureRequirementsTransport2025}.

The final type of transport properties are message properties.
These are set by the application on a per-message basis, and
can be used to control the behavior of individual messages.
These are typically fairly self-explanatory, but do
include certain properties which can override other properties
set for the whole Connection.

For example, the \textit{Ordered}
message property can be used to ensure in-order delivery
of a certain set of messages, even if the connection
does not require it. This naturally requires that the underlying
transport protocol supports switching between ordered and
unordered delivery of messages. This can for example be
supported by an implementaiton of QUIC, but it is not mandatory.

\subsection{Candidate Gathering} \label{sec:candidate-gathering}
When initiating a preconnection, a TAPS system will gather a
list of candidates for establishing a connection. This process
is known as Candidate Gathering. The candidates are based on
the specified local and remote endpoints of the preconnection,
as well as the protocols which are available to use. 

RFC9623 conceptualises these combinations as a tree,
with the first level of child nodes representing the
possible Local Endpoints, which are resolved to possible local IP
addresses when building the tree \cite{brunstromImplementingInterfacesTransport2025}.
The second level of the candidate tree represent the possible protocol
options which can be used. This includes not only the protocol
itself, but also any combination of options for this protocol which
influence connection establishment, such as ALPN values.

The third and final level of child nodes represent any resolved
Remote Endpoints, which have been resolved into specific IP addresses.

Which branches are possible at each level can be restricted
While the primary use of Selection Properties are to restrict
which protocols are selected, they can also be used to affect
other layers in the Candidate Gathering tree.
For example the possible types of interfaces
in the first level can be restricted via the "interface" selection
property and the possible addresses resolved in the third level can
for example be restricted by requiring the use of temporary local
addresses, as these are only featured in IPv6.

While RFC9623 uses a tree to conceptualize the different combinations
of candidates, it places no bearing on the actual implementation
in a real TAPS system, it does however recommend to branch in the
same order as the RFC, to place concepts more important to a
eventual connection higher in the conceptual tree.

After all the candidates have been gathered, a TAPS implementation
has to sort them according to the requirements of the user. The
candidates are sorted from most to least preferred, with the
leaf nodes representing the candidates that are most desirable.

After all the candidates have been gathered and sorted, a single candidate
needs to be selected to be used for this connection. This is done
through Candidate Racing.

\subsection{Candidate Racing} \label{candidateRacing}
When a tree or other representation of the candidates has
been constructed, a single candidate needs to be selected for use
by the final application. This is done through
Candidate Racing. Candidate Racing is the process of establishing
a connection using a single protocol over a single path with
a single endpoint as quickly as possible, optimizing for the
user requirements while minimizing the network load
\cite{brunstromImplementingInterfacesTransport2025}.

This is achieved by trying to establish connections with
the various candidates and finally selecting one of them.
The exact process of racing the candidates and selecting
a winner can vary by implementation. RFC 9623 lists three
ways of performing candidate racing:

\begin{enumerate}
    \item Simultaneous
    \item Staggered
    \item Failover
\end{enumerate}

Simultaneous racing is the process of trying to establish
connections with all candidates at the same time. RFC 9623
discourages this method, as it can lead to an unecessarily large
network load.

Staggered racing is the process of trying to establish connections
with several candidates at once, but with a delay between the start
of each candidate. This delay is configurable, but should be based
on the expected handshake time
\cite{brunstromImplementingInterfacesTransport2025}. 
Since the most desirable candidate will be raced first, it
can be chosen if it connects successfully thereby limiting the
number
of needed attempts to find a connection. When racing IP addresses,
TAPS implementations should follow the algorithm described
in RFC 8305 \cite{schinaziHappyEyeballsVersion2017}. The relevance
of this algorithm for TAPS is described in section \ref{happyEyeballs}.

Failover racing is the process of racing a single candidate
at a time, only attempting the next candidate if the previous
fails. This can be used when an application has a stirng preference
of one branch over another.

Connection establishment has been completed when a single candidate,
has successfully established a connection with a remote endpoint, or
when all candidates have failed. 

\section{Happy Eyeballs} \label{happyEyeballs}
Happy Eyeballs, which is described in RFC 86305 is an algorithm
for racing IPv4 and  IPv6 addresses on dual-stack hosts
\cite{schinaziHappyEyeballsVersion2017}. The Candidate Gathering
and Candidate Racing algorithm found in TAPS is based on Happy Eyeballs,
so they share a lot of central features.
Happy Eyeballs implements a staggered racing algorithm with a preference for
IPv6, but uses interleaved IPv6 and IPv4 addresses in order to maximize
IPv6 usage but minimize latency due to IPv6 brokenness.

The first version of Happy Eyeballs was published in 2012 and describes an algorithm
for racing an IPv4 and IPv6 connection to a server, giving a preferential
head start to the IPv6 connection \cite{wingHappyEyeballsSuccess2012}.
Happy Eyeballs version 2 improves
upon problems within Happy Eyeballs version 1 by also racing the
IPv6 and IPv4 DNS queries, avoiding DNS slow family blocking \cite{schinaziHappyEyeballsVersion2017}.

This concept of racing is a central concept of TAPS, where it is
taken one step further to also races all combinations of local interface,
remote endpoint and underlying transport protocol which fit with the
specifications of the user. RFC 8305 lists four steps for protocol racing,
which conceptually divide into Candidate Gathering and Candidate Racing 
TAPS system.

\begin{enumerate}
  \item Initiation of asynchronous DNS queries
  \item Sorting of resolved destination addresses
  \item Initiation of asynchronous connection attempts
  \item Establishment of one connection, which cancels all other attempts
\end{enumerate} \cite{schinaziHappyEyeballsVersion2017}

Here, step one and two map into the TAPS concept of Candidate gathering
while step two and three map into the TAPS concept of Candidate Racing.
By expanding the racing from only to IP-version to also include transport
protocol, TAPS makes use of techiniques originally used to help IPv6
adoption and user experience, to also help transport protocol adoption.
TAPS cannot only discover that whether a link can be made using IPv6
or only IPv4, but also whether the connection could be made by QUIC
or if it has to made by TCP for example. This will not only lead to
an improved user experience, but also handle fallback-logic which otherwise
would have to be written on a per-application basis.

\section{What is QUIC}
One of the key elements of TAPS is to make new transport protocols
easier to adopt, and one of the most important recent innovations on the
transport layer is QUIC.
QUIC is a protocol developed by Google released to the
public in 2013 described in RFC 9000
\cite{iyengarQUICUDPBasedMultiplexed2021}. QUIC is a transport
layer protocol which is designed to be a replacement for TCP.
QUIC is a connection-oriented protocol which is built on top of
UDP. QUIC enables 0-RTT connection establishment, multistreaming
and integrates encryption as a mandatory feature of the protocol.

While QUIC does present an overarching connection, the
abstraction for each actual data transfer is called a stream
A stream is an either either
unidirectional or biderectional ordered stream of data, much
like what an application sees when using TCP. Streams can be created and
closed by either endpoint, and multple streams can exist over a
single QUIC connection. QUIC mandates support for in-order deliver within a single
stream, but not between streams. This way, a new straem can be
used for each object transferred, solving the issue of
head-of-line blocking which is present in TCP.

The format of a QUIC stream frame as described in RFC 9000 can
be seen in \cref{lst:stream-frame-format}.

\begin{listing}[tbp]
	\centering
	\begin{lstlisting}
        STREAM Frame {
          Type (i) = 0x08..0x0f,
          Stream ID (i),
          [Offset (i)],
          [Length (i)],
          Stream Data (..),
        }
    \end{lstlisting}
	\caption{Format of a QUIC stream frame \cite{iyengarQUICUDPBasedMultiplexed2021}}
	\label{lst:stream-frame-format}
\end{listing}

A given stream can exist in several different states.
These states differ whether the stream is sending or receiving
data. For a stream sending data, the states are:
\texttt{Ready}, \texttt{Send}, \texttt{Data Sent}, \texttt{Reset Sent},
\texttt{Data Recvd}, and \texttt{Reset Recvd}.
A receiving streams states maps the perceived states of the
sending stream. This means that there are some states
which are not represented in the receiving stream, as they
cannot be inferred from the data received. The states of a
receiving stream are: 
\texttt{Recv}, \texttt{Size Known}, \texttt{Data Recvd}, \texttt{Reset Recvd},
\texttt{Data Read}, and \texttt{Reset Read}.
A bidirectional stream consists of both sending and receiving parts,
and can therefore exist in states belonging to both sending and
receiving streams as described previously.

% TODO GOT HERE ? 

A flow chart describing the QUIC stream state machine can be
seen in \cref{fig:stream_states}.

\begin{figure}[tbp]
    \centering
    \begin{lstlisting}
       o
       | Create Stream (Sending)
       | Peer Creates Bidirectional Stream
       v
   +-------+
   | Ready | Send RESET_STREAM
   |       |-----------------------.
   +-------+                       |
       |                           |
       | Send STREAM /             |
       |      STREAM_DATA_BLOCKED  |
       v                           |
   +-------+                       |
   | Send  | Send RESET_STREAM     |
   |       |---------------------->|
   +-------+                       |
       |                           |
       | Send STREAM + FIN         |
       v                           v
   +-------+                   +-------+
   | Data  | Send RESET_STREAM | Reset |
   | Sent  |------------------>| Sent  |
   +-------+                   +-------+
       |                           |
       | Recv All ACKs             | Recv ACK
       v                           v
   +-------+                   +-------+
   | Data  |                   | Reset |
   | Recvd |                   | Recvd |
   +-------+                   +-------+
    \end{lstlisting}
    \caption{States of a QUIC stream sending data \cite{iyengarQUICUDPBasedMultiplexed2021}}
    \label{fig:stream_states}
\end{figure}

Another important abstraction presented by QUIC is the connection.
A QUIC connection is a shared state between a client and a server 
\cite{iyengarQUICUDPBasedMultiplexed2021}.
A connection is established through a handshake, with support
for 0-RTT handshakes as mentioned previously. The connection
is identified by a connection ID, which is selected by each
endpoint's peer. Packets sent to this peer are identified as
belonging to a given connection, by a destination connection ID
field. Connection IDs make addressing in QUIC independent of addressing in lower layer
protocols \cite{iyengarQUICUDPBasedMultiplexed2021}. Each endpoint
has a unique connection ID, making sure that external observers cannot
identify that the two endpoints are communicating with each
other by looking at the connection ID.

Using a connection ID allows the connectoin to survive a change
in the addressing of the underlying protocol \cite{iyengarQUICUDPBasedMultiplexed2021}.
This functionality is enabled by default, but can be disbled by
one of the peers. When an address is changed, each endpoint has to
independently verify that they can communicate using the new address.
Path validation happens using a special type of stream frame
known as \texttt{PATH\_CHALLENGE} and \texttt{PATH\_RESPONSE}
frames. The \texttt{PATH\_CHALLENGE} frame contains unpredictable
data which must be mirrored back in the \texttt{PATH\_RESPONSE} frame.
If the endpoint validation fails, which happens if no matching response
is received within a given time period, the endpoints can continue
communicating using other known addresses. If there are no valid
addresses, the affected endpoint can signal this to the other
peer by using a \texttt{NO\_VIABLE\_PATH} error.

Much like TCP, QUIC also implements congestion control, there are
however some key differences between the algorithms. A key concept
to QUIC's congestion control is the 'packet number space'. Packet
number spaces are \say{the context in which a packet can be processed
and acknowledged} \cite{iyengarQUICUDPBasedMultiplexed2021}. 
This context reflects the level of encryption and there 
are therefore three separate packet number spaces. Initial space,
for the initial handshake packets, handshake space, for packets
exchanged during the cryptographic handshake, and application
data space, for packets exchanged after the cryptographic handshake
has been completed. Separating these spaces allows QUIC to
avoid spurious retransmission from one space affecting another,
and allows QUIC to handle retransmissions more efficiently, for
example by using a different expected acknowledgement delay for packets sent
in the initial space, as these are not expected to be cached
by the peer. Although the packet number spaces are handled
separately when it comes to acknowledgements, they are still
unified when it comes to congestion control and round trip
time measurements.

Another difference between how QUIC and TCP handle acknowledgements
is that QUIC separates the transmission order from the delivery order.
This is done by using stream offset as delivery order, and packet number as
transmission order. This removes ambiguity about whether the
original or the retransmitted packet is acknowledged in the case
of retransmissions, making RTT measurements in these cases more accurate
\cite{iyengarQUICUDPBasedMultiplexed2021}. Another key difference
is that instead of TCP algorithms such as retransmission timeout (RTO)
and tail-loss probe (TLP), QUIC uses probe timeout (PTO).
PTO is an algorithm which is used to detect lost packets. Every
time a packet requiring an acknowledgement is sent, a timer is started. If
the timer expires, QUIC sends a probe packet and the PTO period is
set to twice its previous value. When an acknowledgement is received,
the PTO period is reset to its previous value, as long as the
connection has moved out of the Initial space. This is beneficial
as it can mitigate lost acknowledgements, and does not rely on
duplicate acknowledgements, making it beneficial in several cases
such as tail-loss. \todo{Make sure this is correct}

For QUIC a single connection can exist over several paths and
congestion control is done on a per-path basis. Multiple streams
can be transmitted over a single connection and data from
multiple stream frames can be transmitted in a single packet.
If a packet is lost, only applications using the streams contained
in that packet will be affected. This is a key difference from TCP,
where a single lost packet will block all other streams of data
over the same connection. This is called head-of-line blocking.
This problem is avoided by QUIC, since QUIC only guarantees
in-order delivery of data on a per-stream basis. This leads
to a better utilization of the available bandwidth, when
compared to multiplexing streams over a single TCP connection.

Future network ossification related to QUIC is also avoided
since QUIC mandates TLS 1.3 as part of the protocol.
By encrypting as much of the packet as possible, QUIC is able to
obscure details about the protocol from middleboxes, and
they cannot grow to depend on the contents of the packet.

As mentioned previously, one of many reasons behind the success
of TCP is
the universal support of sockets, which are implemented in kernel space.
There is no equivalent API for QUIC, and with a relatively
advanced implementation, users have to settle for one of the
existing libraries.


Currently, few applications can be built exclusively on top of
QUIC, as it requires support from both the client and the
server. This means that a client wanting to use QUIC, also
has to have support for fallback code, in case the server
it wants to connect to, does not have the required support.
As mentioned previously, making the adoption of new protocols
such as QUIC is one of the key selling points of TAPS. With TAPS,
the client would be able to specify the requirements for the
connection, and TAPS could select QUIC if it was supported, otherwise
falling back to TCP. This would require no change in logic for
the client, and would not require TAPS to be supported both
by the client and the server.

\section{ALPN}
Application layer protocol negotiation is a feature in
TLS where the application layer protocol which will be
used is negotiated between the server and the client.
transport layer negotitates which application layer protocol
will be used between the server and the client. ALPN is featured
in several encryption protocols such as TLS, and is mandatory
in QUIC. ALPN is defined in RFC 7301 and has become necessary
as more and more application level protocls encrypt their
contents with TLS over port 443 \cite{friedlTransportLayerSecurity2014}.
ALPN allows the client and server to negotiate a common protocol
without adding network round trips. The client features
a list of supported protocols in its client hello message
and the server chooses a single protocol and features it in
its server hello message. If the server does not support any
protocols that the client supports, then it responds with an
error notifying the client of this.

The implementation of ALPN of picoquic caused some issues
between picoquic and CTaps, as discussed in \cref{sec:quicImpl}.

\section{SCTP}
Another transport protocol which could be adopted easier
when using TAPS is Stream Control Transmission Protocol (SCTP).
SCTP provides a reliable, message-oriented transport protocol.
SCTP also provides congestion avoidance and flow control like TCP.
It also has support for multihoming, which helps provide a more
reliabale service than traditional TCP connections \cite{leungMultihomedCommunicationSCTP2013}.
By provdiding a message-oriented transport protocol, SCTP provides
a more realistic abstractoin for how applications actually interact
with the network, as mentioned with TAPS. 
However, since SCTP is a transport layer protocol traditionally
built on top of IP, it provides a good case study of how network
ossification can prevent the adoption of new protocols. Even
SCTP/UDP, which is a version of SCTP built on top of UDP, has
struggled with adoption, showing that not only network ossification
prevents the adoption of new protocols, but also that
the design-time binding of an application to a specific
transport protocol also presents a large obstacle to adoption.

Like TCP, SCTP is connection-oriented, but unlike TCP,
\say{the SCTP association is a broader concept than the TCP connection} \cite{leungMultihomedCommunicationSCTP2013}.
SCTP associations can consist of multiple transport addresses,
which consists of a IP address and SCTP port, each of which
can be used to reach the given endpoint. This allows SCTP to
support multihoming, allowing the communication to migrate
in case one connection becomes unavailable.

Originally designed to support telephony signaling, but
adoption in other areas has been limited. Part of
this reason is the lack of support, for example Windows does
not have native SCTP support \cite{hoggWhatStreamControl}.
Another challenge with SCTP adoption is the same as with
many other transport protocols, network ossification.
As a mitigation for this, SCTP can be used over UDP, which
was designed to allow SCTP to be used in environments
where middleboxes are not aware of SCTP, and therefore
do not support it \cite{tuexenUDPEncapsulationStream2013}.

\section{Dependencies}
\subsection{Libuv}
CTaps is reliant on several
external libraries, both for testing and for core functionality. The largest
and most essential of these are libuv \cite{Libuv2025}. Libuv provides the central
event loop which allows CTaps to be asynchronous. It also provides several quality
of life features, such as built in support for asynchronous interactions with TCP
and UDP, through the \texttt{uv\_tcp\_t} and \texttt{uv\_udp\_t} handles respectively.
When adding protocols which libuv does not have native support for, such as QUIC,
it offers a custom work handle \texttt{uv\_work\_t}, which can be used to define
a custom piece of work. This can be used to define both a long-lived QUIC-based
server, or a short interaction with an existing server. Which underlying handle
is actually used in libuv is naturally abstracted away in CTaps, and for the
end-user of the library the interaction is the same no matter what. \todo {very informal lol}

\subsection{Glib}
In order to increase the reliability of internal data structures, as well as
minimizing the amount of time spent reinventing the wheel, CTaps relies on the
Glib library. Glib is a low level core library which is used at the core of
larger libraries such as GTK and GNOME \todo{citation}. Glib provides several
key data structures used internally in CTaps. Among others, it provides the
Hash Map used to keep track of received Connections in a Listener, the tree
used to gather candidates in the Candidate Gathering algorithm, among others.
While adding an external library, especially in core functionaly such as networking,
poses a risk, the benefits were deemed to be greater than the costs. Primarily,
this consideration was made on the basis of Glib being a relatively high-profile
library used in several large projects already, in addition to the benefits of
increased development speed, since CTaps would not have to implement its own 
basic data structures. Using a well-tested external library has the added benefit
of avoiding bugs, as if CTaps rolled its own data structures, it would have
to worry not only about bugs in the library logic, but also in the
data structures themselves.

In order to minimize how much a library consumers would have to
adhere to Glib, it was decided that the external API should be
kept free of Glib data structures, and that all external functions
should only take parameters defined in standard C or in CTaps itself.

\subsection{picoquic} \label{sec:picoquic}
Because the applicability of running QUIC under a TAPS system is one of
the main research questions behind CTaps, supporting QUIC was a high
priority. Because libuv does not feature native QUIC support like it
does for TCP and UDP, an external library was necessary. Because of
the exploratory nature of CTaps, picoquic \cite{christianhuitemaPicoquic} was chosen over other, larger
C QUIC implementations such as MsQuic \cite{MsQuic} from Microsoft or XQUIC \cite{XQUIC} from Alibaba 
because of its relative simplicty.

Picoquic aims to be a minimalist implementation of the QUIC standard
and features support for several extensions, which have not yet
been made available to an Application through CTaps. The prime
functionality of using QUIC as a reliable transport is however
accessible and serves as a good proof of concept for future TAPS
implementations.

Picoquic features two modes of operation. Either an Application can
use the built-in event loop \mono{picoquic_packet_loop} or the
Application can supply its own timer, checking if picoquic needs
to perform any I/O at given intervals. These intervals can be fetched
from picoquic by using the method \mono{picoquic_get_next_wake_delay}.
A more detailed description of the CTaps' QUIC implementation can be
seen in \cref{sec:quicImpl}.

\subsection{log.c}
Log.c is a small logging library which provides a simple API for
formatting log messages of different levels of severity, either to
stdout or to a dedicated logging file \cite{logc}.
